% !TEX root = ../main.tex

\section{Context and Motivation}
\label{sec:intro-context}

Machine learning systems are increasingly deployed in consequential domains such as credit scoring, employment screening, healthcare triage, and public resource allocation. Their widespread use has exposed a growing tension between two intertwined requirements: on one side, the demand for predictive accuracy in large, heterogeneous datasets; on the other, the need for transparency and fairness in how decisions are made. Explainable Artificial Intelligence (XAI) has emerged as the research area tasked with addressing this tension, aiming to make the internal logic of black-box models intelligible to humans without compromising their predictive performance.

Within XAI, a distinction is often drawn between \emph{local} and \emph{global} explanation methods. Local explanations describe why a particular instance received a particular prediction, typically by perturbing features or examining nearby counterfactuals. Global explanations, by contrast, aim to characterize the overall behavior of the classifier, summarizing the recurring patterns that drive its predictions across the entire data distribution. Counterfactual-based methods have become a central tool in the local setting because they express importance in an operational form: a feature is relevant when modifying its value flips the prediction. However, aggregating many local counterfactual explanations into a coherent global picture remains an open problem, particularly in the presence of categorical attributes and protected demographic variables.

The present thesis addresses this problem by proposing a pipeline that combines counterfactual-based feature importance with Association Rule Mining (ARM) to extract a structured, global interpretation of a classifier's decision logic. The motivating application domain is income prediction on U.S.\ census data, where concerns about demographic bias intersect with the interpretability challenges of modern tabular classifiers. The goal is not to build a new predictor, nor to claim legal-grade evidence of discrimination, but to expose the patterns through which the classifier organizes its decisions and to distinguish, in a principled way, patterns that reflect modifiable attributes from patterns that rely on protected demographic features.

\section{Thesis Objectives}
\label{sec:intro-objectives}

The thesis pursues four objectives. The first is to adapt an existing boundary-based counterfactual method for feature importance, originally defined for continuous features, to the fully categorical setting obtained after preprocessing the ACS Income dataset. The second is to design a hierarchical Association Rule Mining procedure that operates on the local explanations produced by the adapted method, distinguishing structural dependencies (which features interact) from semantic content (which specific category values materialize those interactions). The third is to propose a semantic classification of the extracted rules that separates \emph{actionable} rules, involving only modifiable attributes, from \emph{biased} rules, involving protected demographic features. The fourth is to validate the pipeline empirically across multiple geographical benchmarks and classifier families, assessing whether the patterns it surfaces are consistent with the socio-economic structure of the underlying populations.

These objectives are motivated by the observation that macroscopic patterns alone do not provide sufficient explanatory depth (they indicate \emph{which} features matter but not \emph{how}), while microscopic patterns alone are too fragmented to yield a global picture. A hierarchical approach can combine the two levels, using macroscopic rules as a structural anchor for a more targeted microscopic analysis.

\section{Main Contributions}
\label{sec:intro-contributions}

The main contributions of this thesis are the following.

\begin{enumerate}
  \item \textbf{Adaptation of BoCSoR to categorical data.} The BoCSoR method, originally defined for continuous features, is adapted to the fully categorical setting obtained after preprocessing. The adaptation introduces a hybrid encoding (one-hot for nominal features, ordinal rank for ordered ones) and a Manhattan distance formulation that recovers the Hamming distance on the nominal component while preserving an ordered notion of similarity on the ordinal component. Interpolation, which is central in the original method but meaningless in a discrete space, is replaced by the use of $k$ observed opposite-class neighbors, aggregated through a union of per-neighbor relevant feature-category pairs.

  \item \textbf{Hierarchical macroscopic-to-microscopic Association Rule Mining.} A two-level ARM procedure is proposed. The macroscopic level operates on feature names only, identifying structural dependencies such as $\texttt{COW} \Rightarrow \texttt{OCCP}$. For each surviving macroscopic rule, the microscopic level examines the subset of transactions that instantiate it and mines value-level rules such as \texttt{COW=Local-Government-Employee} $\Rightarrow$ \texttt{OCCP=Educational-Instruction-And-Library}. The hierarchy reduces combinatorial explosion and provides a coarse-to-fine reading of the classifier's decision logic.

  \item \textbf{Semantic classification of rules (actionable, biased, other).} A conservative, mutually exclusive classification is introduced. Rules involving only features that an individual can in principle modify (education, class of worker, working hours, occupation) are labeled actionable; rules involving at least one protected demographic feature (sex, race) are labeled biased; all remaining rules are labeled other. This classification does not adjudicate the normative or legal status of the rules; it provides an operational lens for interpreting the pipeline's output.

  \item \textbf{Classifier-agnostic framework with systematic grid search.} The pipeline is validated using two structurally different classifiers, CatBoost and MLP, to separate model-specific artifacts from properties of the underlying data. To avoid dependence on a single pair of thresholds, the ARM stage performs a grid search over support and confidence, producing a full landscape of rule counts.

  \item \textbf{Large-scale empirical study on ACS 2024.} The pipeline is applied to five geographical benchmarks constructed from the 2024 American Community Survey via the Folktables framework: New York, Texas, the U.S.\ Northeast, the U.S.\ South, and the set of all U.S.\ states except Alaska (approximately 1.75 million records in total). The empirical study uncovers a systematic geographic asymmetry: more affluent and urbanized benchmarks (Northeast, New York) are dominated by actionable rules, whereas benchmarks with higher poverty concentration (Texas, South) are dominated by biased rules. The South additionally exhibits a unique bias pattern, $\texttt{RAC1P} \Rightarrow \texttt{POBP}$, that does not emerge as dominant elsewhere. Rules on the higher-income boundary are observed only in New York, and only for broad counterfactual neighborhoods.
\end{enumerate}

\section{Thesis Structure}
\label{sec:intro-structure}

The rest of the thesis is organized as follows. Chapter~\ref{cap:state-of-the-art} situates the contribution in the existing literature on counterfactual explanations, feature importance, Association Rule Mining, and fairness in tabular machine learning. Chapter~\ref{cap:data} describes the construction of the five benchmark datasets, the preprocessing pipeline, and the categorical schema adopted for all downstream analyses. Chapter~\ref{cap:method} formalizes the adapted BoCSoR procedure and the two-level hierarchical ARM, including the semantic classification of rules and the pseudocode of the three algorithms that constitute the pipeline. Chapter~\ref{cap:results} reports the experimental setup, the macroscopic results at three support thresholds, the decay analysis, the class asymmetry, the dominant macroscopic patterns, and the microscopic findings for each geographical scope. Chapter~\ref{cap:conclusions} summarizes the main findings, discusses the limitations of the approach, and outlines directions for future work. The appendices report implementation details, configuration and execution information, output artifacts, the software stack, and supplementary experimental data including classifier accuracy and dataset composition.